%% 第二章

\chapter{大语言模型对齐基础理论与相关技术}
\chaptermark{大模型对齐基础理论与相关技术}

本章梳理后续方法所需的理论与技术基础，内容围绕奖励信号构造和策略优化两条线索展开。第一部分介绍强化学习的基本工具以及PPO、GRPO两种策略优化算法；第二部分讨论RLHF如何用偏好数据训练奖励模型并驱动策略更新；第三部分给出DPO如何把显式奖励建模转化为离线偏好优化；第四部分说明IRL如何从专家示例中学习可随策略变化的奖励函数，并由此引出奖励函数更新的稳定性问题。通过这一组织，本章把大模型偏好对齐问题逐步收束到“奖励信号如何表示、更新与约束”这一主线，为第三章提出基于IRL的PIRO方法做铺垫。

\section{强化学习与策略优化基础}

大模型对齐可以看作在奖励信号引导下调整生成策略的过程。本节先给出强化学习的基本表述，再介绍PPO和GRPO这两种在对齐实验中常用的策略优化算法。

\subsection{马尔可夫决策过程}

强化学习（Reinforcement Learning, RL）是研究智能体如何通过与环境交互学习最优决策策略的机器学习范式\cite{sutton1999policy}。其标准数学框架为马尔可夫决策过程（Markov Decision Process, MDP），由五元组 $\mathcal{M} = (\mathcal{S}, \mathcal{A}, \mathcal{P}, \mathcal{R}, \gamma)$ 定义：$\mathcal{S}$ 为状态空间，$\mathcal{A}$ 为动作空间，$\mathcal{P}: \mathcal{S} \times \mathcal{A} \times \mathcal{S} \to [0,1]$ 为状态转移概率核，对所有 $(s,a)$ 满足 $\sum_{s' \in \mathcal{S}} \mathcal{P}(s'|s,a) = 1$；$\mathcal{R}: \mathcal{S} \times \mathcal{A} \to \mathbb{R}$ 为奖励函数；$\gamma \in [0,1)$ 为折扣因子，权衡即时奖励与未来奖励的相对权重。MDP的核心假设是马尔可夫性质：下一状态 $s_{t+1}$ 的条件分布仅取决于当前状态 $s_t$ 与动作 $a_t$，与历史状态序列无关，即 $\mathcal{P}(s_{t+1}|s_t,a_t,\ldots,s_0,a_0) = \mathcal{P}(s_{t+1}|s_t,a_t)$。

智能体与环境的交互过程如下。在时刻 $t$，智能体观察当前状态 $s_t$，策略 $\pi$ 据此给出各个可选动作的概率分布，智能体再从中采样得到动作 $a_t$；环境根据当前状态和动作转移到下一状态 $s_{t+1}$，并返回即时奖励 $r_t$。强化学习的目标是寻找最优策略 $\pi^*$，使从初始状态出发的期望累积折扣奖励最大：
\begin{equation}
    J(\pi) = \mathbb{E}_{\tau \sim \pi}\left[\sum_{t=0}^{\infty} \gamma^t r_t\right] = \mathbb{E}_{\tau \sim \pi}\left[\sum_{t=0}^{\infty} \gamma^t \mathcal{R}(s_t, a_t)\right]
    \label{eq:rl_objective}
\end{equation}
其中 $\tau = (s_0, a_0, s_1, a_1, \ldots)$ 为策略 $\pi$ 诱导的轨迹，初始状态 $s_0 \sim \rho_0$ 服从初始状态分布。

在大语言模型对齐场景中，文本生成可以理解为一个序贯决策过程。给定输入提示后，模型并非一次性给出完整回答，而是在每个位置依据提示和已生成的上下文预测下一个词元。当前“状态”是提示与已生成文本构成的上下文，“动作”是下一步要生成的词元（token）；模型每生成一个新词元，就将其追加到上下文中作为下一步决策的状态；当生成结束标记（EOS token）或达到长度上限时，本次生成结束。语言模型的词汇表通常包含数万乃至十余万个词元，动作空间很大，生成序列也较长，传统表格式强化学习方法难以适用。RLHF通常直接以语言模型作为参数化策略，通过采样响应估计策略优化目标。在奖励设计上，RLHF常用序列级标量奖励：完整响应生成后由奖励模型对整段回答统一评分，而不是对每一个中间生成步骤单独打分\cite{ouyang2022training}。


\subsection{价值函数与Bellman方程}

价值函数是强化学习理论的基本工具，它把策略的长期性能浓缩为状态或状态-动作对上的标量，便于策略评估与改进。给定策略 $\pi$，状态价值函数 $V^\pi: \mathcal{S} \to \mathbb{R}$ 与动作价值函数（Q函数）$Q^\pi: \mathcal{S} \times \mathcal{A} \to \mathbb{R}$ 分别定义为：
\begin{equation}
    V^\pi(s) = \mathbb{E}_{\tau \sim \pi}\left[\sum_{t=0}^{\infty} \gamma^t r_t \,\bigg|\, s_0 = s\right], \quad
    Q^\pi(s, a) = \mathbb{E}_{\tau \sim \pi}\left[\sum_{t=0}^{\infty} \gamma^t r_t \,\bigg|\, s_0 = s, a_0 = a\right]
    \label{eq:value_functions}
\end{equation}
两者的关系为 $V^\pi(s) = \mathbb{E}_{a \sim \pi(\cdot|s)}[Q^\pi(s,a)] = \sum_{a} \pi(a|s) Q^\pi(s,a)$，以及 $Q^\pi(s,a) = \mathcal{R}(s,a) + \gamma \mathbb{E}_{s' \sim \mathcal{P}(\cdot|s,a)}[V^\pi(s')]$。

Bellman期望方程（Bellman Expectation Equation）刻画了价值函数所满足的递归结构。对于策略 $\pi$，$V^\pi$ 与 $Q^\pi$ 分别满足：
\begin{align}
    V^\pi(s) &= \sum_{a} \pi(a|s) \left[\mathcal{R}(s,a) + \gamma \sum_{s'} \mathcal{P}(s'|s,a) V^\pi(s')\right]
    \label{eq:bellman_v} \\
    Q^\pi(s,a) &= \mathcal{R}(s,a) + \gamma \sum_{s'} \mathcal{P}(s'|s,a) \sum_{a'} \pi(a'|s') Q^\pi(s',a')
    \label{eq:bellman_q}
\end{align}
式~\eqref{eq:bellman_v}与式~\eqref{eq:bellman_q}分别构成关于 $V^\pi$ 和 $Q^\pi$ 的线性方程组，在有限状态空间中可直接求解（策略评估），也可作为时序差分学习（Temporal Difference Learning）的更新目标。借助算子理论，Bellman期望算子 $\mathcal{T}^\pi$ 在无穷范数下以 $\gamma$ 为收缩系数，且 $V^\pi$ 是 $\mathcal{T}^\pi$ 的唯一不动点，由此保证策略评估的收敛性。

Bellman最优方程（Bellman Optimality Equation）则刻画了最优价值函数 $V^* = \max_\pi V^\pi$ 与 $Q^* = \max_\pi Q^\pi$ 所满足的非线性递归结构：
\begin{align}
    V^*(s) &= \max_{a} \left[\mathcal{R}(s,a) + \gamma \sum_{s'} \mathcal{P}(s'|s,a) V^*(s')\right]
    \label{eq:bellman_opt_v} \\
    Q^*(s,a) &= \mathcal{R}(s,a) + \gamma \sum_{s'} \mathcal{P}(s'|s,a) \max_{a'} Q^*(s',a')
    \label{eq:bellman_opt_q}
\end{align}
最优价值函数 $V^*$ 是最优Bellman算子 $\mathcal{T}^*$ 的唯一不动点，$\mathcal{T}^*$ 同样以 $\gamma$ 为收缩系数。经典Q-learning算法\cite{watkins1992q}以式~\eqref{eq:bellman_opt_q}为更新目标，通过随机梯度下降逼近 $Q^*$，进而间接获得最优策略 $\pi^*(s) = \arg\max_a Q^*(s,a)$。

优势函数（Advantage Function）$A^\pi(s, a)$ 定义为动作价值函数与状态价值函数之差：
\begin{equation}
    A^\pi(s, a) = Q^\pi(s, a) - V^\pi(s)
    \label{eq:advantage}
\end{equation}
$A^\pi(s,a)$ 度量在状态 $s$ 下选择动作 $a$ 相对于策略均值行为的收益增量：$A^\pi(s,a) > 0$ 时动作 $a$ 优于均值水平，$A^\pi(s,a) < 0$ 时该动作低于均值水平。优势函数在策略梯度估计中作用直接：用 $A^\pi$ 代替 $Q^\pi$ 作为梯度权重，可在不改变梯度期望的前提下大幅降低估计方差，这是广义优势估计（Generalized Advantage Estimation, GAE）的基本思想，下一小节将详细讨论。

\subsection{策略梯度方法}

策略梯度方法直接对策略参数进行优化，适用于大规模连续动作空间或高维离散动作空间（如语言模型的词元空间）。策略梯度定理\cite{sutton1999policy}给出了目标函数 $J(\pi_\phi) = \mathbb{E}_{s_0 \sim \rho_0}[V^{\pi_\phi}(s_0)]$ 关于参数 $\phi$ 的精确梯度表达式：
\begin{equation}
    \nabla_\phi J(\pi_\phi) = \mathbb{E}_{s \sim d^{\pi_\phi}, a \sim \pi_\phi(\cdot|s)}\left[\nabla_\phi \log \pi_\phi(a|s) \cdot Q^{\pi_\phi}(s, a)\right]
    \label{eq:policy_gradient}
\end{equation}
其中 $d^{\pi_\phi}(s) = (1-\gamma)\sum_{t=0}^{\infty}\gamma^t \mathbb{P}_{\pi_\phi}(s_t=s)$ 为策略 $\pi_\phi$ 诱导的折扣状态占用度量，$\mathbb{P}_{\pi_\phi}(s_t=s)$ 表示在执行策略 $\pi_\phi$ 时第 $t$ 步访问状态 $s$ 的概率。式~\eqref{eq:policy_gradient}把策略梯度写为对当前策略 $\pi_\phi$ 的期望，可通过蒙特卡洛采样得到无偏估计，无需已知环境动力学 $\mathcal{P}$。

最基本的策略梯度算法REINFORCE通过蒙特卡洛方法估计式~\eqref{eq:policy_gradient}：以实际累积折扣奖励 $G_t = \sum_{k=t}^T \gamma^{k-t} r_k$ 作为 $Q^{\pi_\phi}(s_t, a_t)$ 的估计，所得梯度无偏但方差较大；引入不依赖动作的基线函数 $b(s_t)$（通常取 $V^{\pi_\phi}(s_t)$ 的估计）能在保持无偏性的同时降低方差。

演员-评论员（Actor-Critic）框架同时维护策略网络（演员，Actor）与价值网络（评论员，Critic）：Actor以Critic提供的优势函数估计 $\hat{A}(s_t, a_t) = r_t + \gamma V_\psi(s_{t+1}) - V_\psi(s_t)$ 作为策略梯度权重；Critic通过最小化TD误差在线更新，相较纯REINFORCE方差更低。

广义优势估计（Generalized Advantage Estimation, GAE）\cite{schulman2017proximal}给出了一种可在偏差与方差之间灵活权衡的优势函数估计方案。定义TD残差 $\delta_t = r_t + \gamma V_\psi(s_{t+1}) - V_\psi(s_t)$，GAE优势估计为TD残差的指数加权和：
\begin{equation}
    \hat{A}_t^{\text{GAE}(\gamma,\lambda)} = \sum_{l=0}^{\infty} (\gamma\lambda)^l \delta_{t+l}
    \label{eq:gae}
\end{equation}
其中 $\lambda \in [0,1]$ 是控制偏差-方差权衡的超参数：$\lambda = 0$ 时退化为单步TD残差 $\hat{A}_t^{\text{GAE}} = \delta_t$，高偏差低方差；$\lambda = 1$ 时退化为蒙特卡洛优势估计 $\hat{A}_t^{\text{GAE}} = \sum_l \gamma^l \delta_{t+l} = G_t - V_\psi(s_t)$，低偏差高方差。工程实现中 $\lambda$ 一般取 $[0.9, 0.99]$，在偏差与方差之间折中。GAE是现代策略梯度方法的标准组件，在RLHF框架下的PPO实现中被广泛使用\cite{zheng2023secrets}。

\subsection{近端策略优化}

近端策略优化（Proximal Policy Optimization, PPO）\cite{schulman2017proximal}是当前常用的策略梯度算法，也是大语言模型RLHF后训练的主流优化方法\cite{ouyang2022training}。PPO在重要性采样框架下做离线策略更新，借助截断机制约束每次更新的幅度，在样本效率和训练稳定性之间取得平衡。

PPO的基础是替代目标函数（Surrogate Objective）。设 $\pi_{\phi_{\text{old}}}$ 为采样时所用的旧策略，$\pi_\phi$ 为当前待优化策略，重要性采样比率定义为：
\begin{equation}
    \rho_t(\phi) = \frac{\pi_\phi(a_t|s_t)}{\pi_{\phi_{\text{old}}}(a_t|s_t)}
    \label{eq:importance_ratio}
\end{equation}
无截断替代目标 $\mathcal{L}^{\text{CPI}}(\phi) = \mathbb{E}_t[\rho_t(\phi) \hat{A}_t]$ 在 $\pi_\phi = \pi_{\phi_{\text{old}}}$ 时与策略梯度目标等价，但对大幅更新没有限制。PPO-Clip通过对比率截断处理这一问题：
\begin{equation}
    \mathcal{L}^{\text{CLIP}}(\phi) = \mathbb{E}_t\left[\min\left(\rho_t(\phi) \hat{A}_t,\ \text{clip}(\rho_t(\phi), 1-\epsilon, 1+\epsilon) \hat{A}_t\right)\right]
    \label{eq:ppo_clip}
\end{equation}
其中超参数 $\epsilon \in (0.1, 0.3)$ 控制截断区间宽度：当 $\hat{A}_t > 0$ 时截断上界 $1+\epsilon$ 阻止过度强化，当 $\hat{A}_t < 0$ 时截断下界 $1-\epsilon$ 防止过度惩罚，取二者最小值使目标函数对策略变化保持保守。PPO实际训练中除策略目标外还会同时训练价值函数并加入探索正则项。若以最大化目标的形式写出，其完整形式为：
\begin{equation}
    \mathcal{L}^{\text{PPO}} = \mathcal{L}^{\text{CLIP}} - c_v \mathcal{L}^{\text{value}} + c_e \mathcal{H}(\pi_\phi)
\end{equation}
其中 $\mathcal{L}^{\text{value}}$ 为价值函数预测值与实际回报之间的均方误差，用于训练评论员网络；$\mathcal{H}(\pi_\phi)$ 为策略熵，用于让策略保持一定随机性；$c_v$ 和 $c_e$ 分别控制价值函数损失与熵正则项在总目标中的权重。语言模型RLHF实现中，策略熵正则与针对参考策略的KL散度惩罚通常一并使用：前者让策略保留一定探索性，后者限制策略偏离参考模型过远，使奖励分数与语言能力达到折中\cite{zheng2023secrets}。PPO-KL变体则用自适应KL惩罚替代截断，根据实际KL散度与目标KL散度的比较动态调节惩罚系数，更适合需要精确控制策略偏移的场景。

大语言模型RLHF的实际实现中，PPO训练存在若干工程难点\cite{zheng2023secrets}：价值网络须与策略网络同步维护，占用大量GPU显存；要等到完整响应生成后才能计算奖励，样本效率偏低；多项超参数（学习率、截断系数 $\epsilon$、GAE参数 $\lambda$、KL系数、批次大小等）的联合调节较为复杂，对最终对齐效果影响明显。这些问题催生了更轻量的策略优化方法，GRPO是近年来在大模型后训练中常用的代表性方法。

\subsection{组相对策略优化}

组相对策略优化（Group Relative Policy Optimization, GRPO）\cite{shao2024deepseekmath}可以看作PPO在大语言模型场景下的简化版本。PPO通常需要额外训练价值函数来估计每个样本的优势值，GRPO则直接在同一提示下采样多条响应，用组内奖励的相对高低构造优势估计，减少对价值网络的依赖。

具体而言，对同一提示 $x$，从当前策略中采样一组响应 $\{y_1, y_2, \ldots, y_G\}$，再由奖励模型或任务评分函数得到对应奖励 $\{r_1, r_2, \ldots, r_G\}$。GRPO以组内归一化奖励作为优势估计：
\begin{equation}
    \hat{A}_i = \frac{r_i - \mathrm{mean}(\{r_j\}_{j=1}^{G})}{\mathrm{std}(\{r_j\}_{j=1}^{G}) + \epsilon}, \quad i=1,2,\ldots,G,
    \label{eq:grpo_advantage}
\end{equation}
其中，$G$ 为同一提示下采样的响应数量，$r_i$ 为第 $i$ 条响应的奖励，$\epsilon$ 为防止分母为零的小常数。该优势估计的含义直接：同一组中得分高于平均水平的响应被强化，低于平均水平的响应被压低。

在策略更新时，GRPO仍沿用PPO的概率比截断思想，只是把优势估计替换为组内相对优势。其目标可写为：
\begin{equation}
    \mathcal{L}^{\text{GRPO}}(\phi)=\mathbb{E}_{x,\{y_i\}_{i=1}^{G}}\left[\frac{1}{G}\sum_{i=1}^{G}\min\left(\rho_i(\phi)\hat{A}_i,\ \text{clip}(\rho_i(\phi),1-\epsilon_c,1+\epsilon_c)\hat{A}_i\right)\right],
    \label{eq:grpo_objective}
\end{equation}
其中，$\rho_i(\phi)=\pi_\phi(y_i|x)/\pi_{\phi_{\mathrm{old}}}(y_i|x)$ 为新旧策略在第 $i$ 条响应上的概率比，$\epsilon_c$ 为截断系数。实际对齐训练中，GRPO同样会结合参考策略KL惩罚，限制模型相对SFT模型的偏离幅度。

相较PPO，GRPO的主要优势在于省去独立价值函数，降低显存占用和实现复杂度；组内相对比较也更符合语言模型一次采样多条候选回答再比较质量的训练方式。但GRPO并没有改变RLHF的基本结构：策略优化仍依赖外部奖励模型或评分函数。如果奖励信号本身存在偏差或在策略分布变化后失效，GRPO仍会继承奖励模型过优化问题。因此，PPO和GRPO主要回答“给定奖励后如何更新策略”，而无法单独解决“奖励是否可靠”的问题。

图~\ref{fig:ppo_grpo_advantage}进一步对比了PPO与GRPO在优势估计环节上的差异。二者都保留概率比截断来限制策略更新幅度，但PPO依赖价值函数和GAE估计优势，GRPO直接用同一提示下多条响应的组内相对奖励构造优势，因此后者在大模型后训练中具有更低的工程开销。

\begin{figure}[htbp]
    \centering
    \resizebox{0.98\textwidth}{!}{%
    \begin{tikzpicture}[
        font=\small,
        stage/.style={draw, rounded corners, thick, fill=gray!7, minimum width=16.9cm, minimum height=1.35cm},
        labelbox/.style={font=\normalsize\bfseries, align=left, text width=1.8cm},
        box/.style={draw, rounded corners, thick, align=center, minimum width=2.35cm, minimum height=0.82cm, text width=2.35cm},
        arrow/.style={->, thick}
    ]
        \node[stage] at (8.35,1.25) {};
        \node[labelbox] at (1.05,1.25) {PPO\\优势估计};
        \node[box, fill=blue!8] (ppo_sample) at (2.9,1.25) {采样响应\\$(x,y)$};
        \node[box, fill=blue!8] (ppo_reward) at (5.9,1.25) {奖励评分\\$r_\psi(x,y)$};
        \node[box, fill=blue!8] (ppo_value) at (8.9,1.25) {价值函数\\$V_\omega$};
        \node[box, fill=blue!8] (ppo_adv) at (11.9,1.25) {GAE优势\\$\hat{A}_t$};
        \node[box, fill=blue!8] (ppo_update) at (14.9,1.25) {PPO-Clip\\更新策略};
        \draw[arrow] (ppo_sample) -- (ppo_reward);
        \draw[arrow] (ppo_reward) -- (ppo_value);
        \draw[arrow] (ppo_value) -- (ppo_adv);
        \draw[arrow] (ppo_adv) -- (ppo_update);

        \node[stage] at (8.35,-1.25) {};
        \node[labelbox] at (1.05,-1.25) {GRPO\\优势估计};
        \node[box, fill=green!10] (grpo_sample) at (2.9,-1.25) {同一提示\\采样$G$条};
        \node[box, fill=green!10] (grpo_reward) at (5.9,-1.25) {奖励评分\\$r_1,\ldots,r_G$};
        \node[box, fill=green!10] (grpo_stat) at (8.9,-1.25) {组内统计\\均值/标准差};
        \node[box, fill=green!10] (grpo_adv) at (11.9,-1.25) {相对优势\\$\hat{A}_i$};
        \node[box, fill=green!10] (grpo_update) at (14.9,-1.25) {GRPO-Clip\\更新策略};
        \draw[arrow] (grpo_sample) -- (grpo_reward);
        \draw[arrow] (grpo_reward) -- (grpo_stat);
        \draw[arrow] (grpo_stat) -- (grpo_adv);
        \draw[arrow] (grpo_adv) -- (grpo_update);
    \end{tikzpicture}}
    \caption{PPO与GRPO优势估计机制对比}
    \label{fig:ppo_grpo_advantage}
\end{figure}

\subsection{KL约束策略优化与奖励—策略等价关系}

在大语言模型对齐中，若直接以奖励模型评分作为唯一优化目标，策略容易偏离原始语言模型并产生退化输出。因此，RLHF在奖励模型评分之外通常加入对参考策略的约束；这一带约束的策略优化形式也是后续DPO理论推导的基础。单纯追求更高的奖励模型分数往往导致策略过度优化（over-optimization）：模型可能生成奖励模型评分极高、但人类看来质量很低的回答，并逐渐损害预训练阶段获得的语言能力。标准做法是在优化目标中引入对参考策略 $\pi_{\text{ref}}$（通常为SFT后的模型）的KL散度惩罚项，得到带约束的强化学习目标\cite{ouyang2022training}：
\begin{equation}
    \max_{\pi}\, \mathbb{E}_{x \sim \mathcal{D},\, y \sim \pi(\cdot|x)}\left[r(x, y)\right] - \beta \cdot D_{\text{KL}}\left(\pi(\cdot|x) \,\|\, \pi_{\text{ref}}(\cdot|x)\right)
    \label{eq:kl_constrained}
\end{equation}
其中 $\beta > 0$ 为KL惩罚系数，$D_{\text{KL}}(\pi \| \pi_{\text{ref}}) = \mathbb{E}_{y \sim \pi}[\log \pi(y|x) - \log \pi_{\text{ref}}(y|x)]$。调节 $\beta$ 可以权衡奖励最大化与参考策略相似性：$\beta$ 越大策略越保守，越靠近 $\pi_{\text{ref}}$；$\beta$ 越小策略越激进，越接近纯奖励最大化。

式~\eqref{eq:kl_constrained}存在可解析求解的闭式最优解，这是它的重要理论性质。把它改写为如下等价形式：
\begin{align}
    \max_{\pi}\, \mathbb{E}_{x \sim \mathcal{D}}\left[\mathbb{E}_{y \sim \pi(\cdot|x)}\left[r(x, y) - \beta \log \frac{\pi(y|x)}{\pi_{\text{ref}}(y|x)}\right]\right]
    \label{eq:kl_constrained_v2}
\end{align}
对固定的 $x$，内层关于 $\pi(\cdot|x)$ 的优化是带归一化约束 $\sum_y \pi(y|x) = 1$ 的凹函数最大化问题，由KKT条件可得最优解满足：
\begin{equation}
    \pi^*(y|x) \propto \pi_{\text{ref}}(y|x) \exp\left(\frac{r(x, y)}{\beta}\right)
    \label{eq:boltzmann_unnorm}
\end{equation}
经归一化后的最优策略为：
\begin{equation}
    \pi^*(y|x) = \frac{\pi_{\text{ref}}(y|x) \exp\left(r(x, y)/\beta\right)}{Z(x)}
    \label{eq:boltzmann}
\end{equation}
其中归一化配分函数 $Z(x) = \sum_y \pi_{\text{ref}}(y|x) \exp(r(x,y)/\beta)$ 保证 $\pi^*(\cdot|x)$ 的概率归一性。该Boltzmann分布与统计物理存在类比：把 $-r(x,y)/\beta$ 视为能量、$\beta$ 视为逆温度，最优策略恰好对应该能量函数的Gibbs分布。

以下给出式~\eqref{eq:boltzmann}的完整推导。对于固定的 $x$，设Lagrange函数为：
\begin{equation}
    \mathcal{L}(\pi, \lambda) = \sum_y \pi(y|x)\left[r(x,y) - \beta \log \frac{\pi(y|x)}{\pi_{\text{ref}}(y|x)}\right] - \lambda\left(\sum_y \pi(y|x) - 1\right)
\end{equation}
对 $\pi(y|x)$ 求偏导并令其为零：
\begin{equation}
    \frac{\partial \mathcal{L}}{\partial \pi(y|x)} = r(x,y) - \beta \log \frac{\pi(y|x)}{\pi_{\text{ref}}(y|x)} - \beta - \lambda = 0
\end{equation}
由此解出 $\pi(y|x) = \pi_{\text{ref}}(y|x) \exp\left(\frac{r(x,y) - \beta - \lambda}{\beta}\right)$，代入归一化条件 $\sum_y \pi(y|x) = 1$ 确定 $\lambda$，化简后得到式~\eqref{eq:boltzmann}。由于目标函数关于 $\pi(\cdot|x)$ 严格凹（负熵项严格凸），该解为唯一全局最优解。

对式~\eqref{eq:boltzmann}两边取对数并整理，可得到奖励函数与最优策略之间的等价关系：
\begin{equation}
    r(x, y) = \beta \log \frac{\pi^*(y|x)}{\pi_{\text{ref}}(y|x)} + \beta \log Z(x)
    \label{eq:reward_reparam}
\end{equation}
式~\eqref{eq:reward_reparam}给出了一个等价关系：在KL约束RL目标下，奖励函数与最优策略一一对应，附加项 $\beta \log Z(x)$ 仅依赖于提示 $x$ 而与响应 $y$ 无关。这一性质是DPO推导的数学基础，第\ref{sec:dpo}节会详细给出。

\section{基于人类反馈的强化学习}
\label{sec:rlhf}

\subsection{人类偏好数据的收集与处理}

基于人类反馈的强化学习（Reinforcement Learning from Human Feedback, RLHF）\cite{christiano2017deep,ouyang2022training,ziegler2019fine}是当前产业界部署最广的大模型对齐范式，其基本思想是用人类偏好判断作为监督信号，引导语言模型行为向符合人类价值观的方向调整。RLHF的完整流程包括三个阶段：监督微调（Supervised Fine-Tuning, SFT）、奖励模型训练（Reward Model Training, RM）以及基于PPO或GRPO的策略优化。

图~\ref{fig:rlhf_pipeline}给出了RLHF在大语言模型后训练中的典型三阶段流程。该流程先用人工示例得到可遵循指令的SFT模型，再用成对偏好数据训练奖励模型，最后以奖励模型为代理目标并结合KL约束执行策略优化。

\begin{figure}[htbp]
    \centering
    \resizebox{0.98\textwidth}{!}{%
    \begin{tikzpicture}[
        font=\small,
        stage/.style={draw, rounded corners, thick, fill=gray!7, minimum width=14.8cm, minimum height=1.28cm},
        labelbox/.style={font=\bfseries, align=left, text width=2.2cm},
        box/.style={draw, rounded corners, thick, align=center, minimum width=2.45cm, minimum height=0.82cm, text width=2.45cm},
        arrow/.style={->, thick}
    ]
        \node[stage] at (7.0,2.8) {};
        \node[labelbox] at (0.8,2.8) {阶段一\\监督微调};
        \node[box, fill=green!10] (sft_data) at (3.1,2.8) {人工示例\\$(x,y)$};
        \node[box, fill=green!10] (sft_train) at (6.3,2.8) {监督学习\\最大似然};
        \node[box, fill=green!10] (sft_model) at (9.5,2.8) {SFT模型\\$\pi_{\mathrm{ref}}$};
        \draw[arrow] (sft_data) -- (sft_train);
        \draw[arrow] (sft_train) -- (sft_model);

        \node[stage] at (7.0,0.7) {};
        \node[labelbox] at (0.8,0.7) {阶段二\\奖励建模};
        \node[box, fill=yellow!15] (rm_prompt) at (3.1,0.7) {提示采样\\$x$};
        \node[box, fill=yellow!15] (rm_resp) at (6.3,0.7) {SFT采样响应\\$y_1,\ldots,y_K$};
        \node[box, fill=yellow!15] (rm_pref) at (9.5,0.7) {人工排序\\$y_w \succ y_l$};
        \node[box, fill=yellow!15] (rm_model) at (12.7,0.7) {奖励模型\\$r_\psi(x,y)$};
        \draw[arrow] (rm_prompt) -- (rm_resp);
        \draw[arrow] (rm_resp) -- (rm_pref);
        \draw[arrow] (rm_pref) -- (rm_model);

        \node[stage] at (7.0,-1.4) {};
        \node[labelbox] at (0.8,-1.4) {阶段三\\策略优化};
        \node[box, fill=blue!8] (ppo_sample) at (3.1,-1.4) {当前策略\\采样响应};
        \node[box, fill=blue!8] (ppo_reward) at (6.3,-1.4) {RM评分\\+参考KL};
        \node[box, fill=blue!8] (ppo_update) at (9.5,-1.4) {PPO/GRPO\\更新策略};
        \node[box, fill=blue!8] (ppo_model) at (12.7,-1.4) {对齐后模型\\$\pi_\phi$};
        \draw[arrow] (ppo_sample) -- (ppo_reward);
        \draw[arrow] (ppo_reward) -- (ppo_update);
        \draw[arrow] (ppo_update) -- (ppo_model);
    \end{tikzpicture}}
    \caption{基于人类反馈的强化学习三阶段训练流程}
    \label{fig:rlhf_pipeline}
\end{figure}

偏好数据收集是RLHF流程的起点，其质量直接决定对齐效果的上限。标准流程为：从提示数据库 $\mathcal{D}_{\text{prompt}}$ 抽取提示 $x$；从当前语言模型（或多个版本）采样若干响应 $\{y_1, y_2, \ldots, y_K\}$，通常取 $K=2$ 进行成对比较；由经过培训的标注员按照预定义评估准则（有用性、无害性、诚实性等）比较响应，记录偏好关系 $y_w \succ y_l | x$（$y_w$ 为偏好响应，$y_l$ 为非偏好响应）；最终汇总为偏好数据集 $\mathcal{D} = \{(x^{(i)}, y_w^{(i)}, y_l^{(i)})\}_{i=1}^N$\cite{ouyang2022training}。

标注质量是偏好数据收集的主要难点\cite{cui2023ultrafeedback}。偏好判断带有主观性，标注一致性通常仅为60\%至80\%。常见的提质做法包括多人标注多数表决、黄金标准样本校准、置信度加权训练以及AI辅助标注（如Constitutional AI\cite{bai2022constitutional}）等。

偏好数据的多样性同样关键\cite{cui2023ultrafeedback}。UltraFeedback\cite{cui2023ultrafeedback}聚合多种来源（代码、数学、写作、问答等），用GPT-4对响应打分，构建了大规模多维偏好数据集，奖励模型的泛化能力随之改善。奖励差距过小的样本学习信号微弱，差距过大的样本训练价值有限，合理的差距分布是数据构建中需要考虑的因素。

\subsection{Bradley-Terry偏好模型}

RLHF奖励建模的关键，是把人类偏好判断转化为可微分的训练目标，这需要为偏好数据建立概率模型。Bradley-Terry（BT）模型\cite{bradley1952rank}是配对比较数据建模的经典框架，最初用于体育赛事中球队实力的排名估计，后来在信息检索、推荐系统和人类反馈建模中也有应用。

BT模型为每条响应 $y$ 打出一个分数（即奖励值）$r(x,y) \in \mathbb{R}$。给定两条响应 $(y_1, y_2)$，人类偏好 $y_1$ 优于 $y_2$ 的概率为：
\begin{equation}
    \mathbb{P}(y_1 \succ y_2 | x) = \sigma\left(r(x, y_1) - r(x, y_2)\right) = \frac{\exp(r(x,y_1))}{\exp(r(x,y_1)) + \exp(r(x,y_2))}
    \label{eq:bt_model}
\end{equation}
其中 $\sigma(z) = 1/(1+e^{-z})$ 为sigmoid函数。BT模型有几条性质值得注意：（1）对称性：$\mathbb{P}(y_1 \succ y_2 | x) + \mathbb{P}(y_2 \succ y_1 | x) = 1$；（2）传递性弱约束：若 $r(x,y_1) > r(x,y_2) > r(x,y_3)$，则 $\mathbb{P}(y_1 \succ y_3 | x) > \max(\mathbb{P}(y_1 \succ y_2 | x), \mathbb{P}(y_2 \succ y_3 | x))$；（3）奖励平移不变性：奖励函数加任意常数 $c$ 不影响任何配对比较的概率。

BT模型的一个推广是Plackett-Luce模型，把配对比较扩展到 $K$ 候选的完整排名，一次排名任务就能产生 $\binom{K}{2}$ 对比较信息。BT模型隐含一个较强的假设：人类偏好应有相对一致的排序关系。例如回答A优于B、B优于C，那么通常也认为A优于C。实际标注中，人类判断会受语境、主观标准和噪声影响，出现不一致，BT模型也由此产生偏差。即便如此，由于其形式简单、易于优化、实践效果稳定，BT仍是RLHF奖励建模中最常用的偏好建模方法\cite{ouyang2022training}。

\subsection{奖励模型训练}

给定偏好数据集 $\mathcal{D} = \{(x^{(i)}, y_w^{(i)}, y_l^{(i)})\}_{i=1}^N$ 与BT偏好模型，奖励模型 $r_\psi(x,y)$（参数为 $\psi$）通过最大化偏好数据的对数似然进行训练，等价于最小化如下交叉熵损失：
\begin{equation}
    \mathcal{L}_{\text{RM}}(\psi) = -\mathbb{E}_{(x, y_w, y_l) \sim \mathcal{D}}\left[\log \sigma\left(r_\psi(x, y_w) - r_\psi(x, y_l)\right)\right]
    \label{eq:rm_loss}
\end{equation}
该损失函数要求奖励模型对偏好响应 $y_w$ 的评分明显高于非偏好响应 $y_l$：当 $r_\psi(x,y_w) - r_\psi(x,y_l) \gg 0$ 时损失趋近零，差值为负时损失大幅升高。

奖励模型通常以大规模预训练语言模型（如LLaMA-3\cite{dubey2024llama3}、Qwen2.5等）作为初始化，在最后一层Transformer隐状态之上加一个线性映射头，把序列末尾词元的隐状态映射为标量奖励值 $r_\psi(x,y) \in \mathbb{R}$\cite{ziegler2019fine}。初始化策略对训练效果影响明显：以SFT模型（而非预训练基础模型）作为初始化通常能获得更高的偏好预测准确率，因为SFT模型已经具备理解对话内容和遵从指令的基本能力\cite{ouyang2022training}。

奖励模型的主要风险在于：它可能学到偏好数据中的表面规律，而不是真正反映回答质量的信号\cite{gao2023scaling}。例如，若训练数据中较长的回答更常被标为优选，奖励模型就可能形成长度偏好（longer-is-better bias），即便内容质量没有明显提升，也会给更长的回答更高评分\cite{zheng2023secrets}。这种偏差在单独评估奖励模型时未必明显，但在PPO等策略优化过程中会被不断放大：策略会倾向于生成奖励模型喜欢、人类未必真正偏好的回答。因此，奖励模型的可靠性不能只看保留测试集上的偏好预测准确率，还要看它能否对当前策略生成的新样本给出稳定可信的评分，这也是后续分布偏移与奖励过优化问题的来源\cite{gao2023scaling,li2023reward}。

\subsection{基于PPO/GRPO的策略优化}

奖励模型训练完成后，RLHF第三阶段以固定的奖励模型 $r_\psi$ 作为奖励信号，用PPO或GRPO等算法优化语言模型策略 $\pi_\phi$。两类算法的差别主要在优势估计和工程实现上，基本目标相同：提高奖励模型评分，同时限制策略偏离参考模型。标准RLHF的策略优化目标为\cite{ouyang2022training}：
\begin{equation}
    \max_\phi\, \mathcal{J}_{\text{RLHF}}(\phi) = \mathbb{E}_{x \sim \mathcal{D},\, y \sim \pi_\phi(\cdot|x)}\left[r_\psi(x, y) - \beta \cdot \log\frac{\pi_\phi(y|x)}{\pi_{\text{ref}}(y|x)}\right]
    \label{eq:rlhf_obj}
\end{equation}
其中 $\pi_{\text{ref}}$ 为参考策略（SFT模型），KL散度项 $\beta \log(\pi_\phi/\pi_{\text{ref}})$ 在词元级别计算——对生成序列中每个词元的对数概率比求和——以约束策略相对参考模型的偏离幅度。

在实现层面，式~\eqref{eq:rlhf_obj}中的KL项通常被分解为词元级即时奖励：在序列每个位置 $t$，奖励为 $r_t = -\beta [\log \pi_\phi(a_t|s_t) - \log \pi_{\text{ref}}(a_t|s_t)]$，序列终止位置再叠加奖励模型评分 $r_\psi(x,y)$。这样PPO的优势估计和GAE就能直接用于词元级别的时序决策\cite{zheng2023secrets}。GRPO则在同一提示下采样多条响应，用组内相对奖励构造优势估计，减少对价值函数的依赖。

以PPO为例，RLHF的完整训练循环依次为：采样提示、生成响应、计算奖励（含词元级KL惩罚）、GAE优势估计、PPO截断目标梯度更新策略、价值网络更新，循环执行直至收敛\cite{zheng2023secrets}。GRPO中价值函数更新被组内相对优势估计替代，训练链路更轻量。

KL系数 $\beta$ 的调度对训练稳定性影响较大：实际KL散度超过目标值1.5倍时增大 $\beta$，低于目标值的 $1/1.5$ 时减小 $\beta$，使优化过程不至于过保守或过激进。

\subsection{分布偏移与奖励过优化}

RLHF范式存在一个结构性的缺陷：奖励模型在偏好数据上训练完成后参数即被冻结，不会随策略演化更新；策略在PPO或GRPO优化中持续生成新样本，其生成分布会逐渐偏离奖励模型的训练分布。这一现象称为分布偏移（distribution shift），也是RLHF中奖励过优化（reward over-optimization）的根本成因\cite{gao2023scaling}。

Gao等人\cite{gao2023scaling}用大量实验量化了奖励过优化现象：随着策略与参考模型之间的KL散度 $d$ 增大，代理奖励单调递增，而黄金奖励先升后降，近似满足 $r_{\text{gold}}(d) \approx \alpha \sqrt{d} - \beta d$，存在峰值点 $d^* = (\alpha/2\beta)^2$，超过该阈值后策略进入奖励黑客（reward hacking）区域。原因在于：奖励模型在训练分布之外存在系统性高估误差，策略优化会发现并利用这些高估区域，产生代理奖励高、实际质量差的响应。

针对上述问题，常见的缓解方案包括\cite{moskovitz2023fool,gao2023scaling}：加强KL散度约束以压制策略偏移；奖励模型集成以减少单一模型漏洞；在策略优化过程中动态收集当前策略样本的偏好标注并持续更新奖励模型。最后一种方案把奖励模型由静态组件转为动态组件，是缓解分布偏移的更根本路径。RLHF的主要矛盾不只是PPO还是GRPO的选择问题，更在于奖励模型能否在策略持续变化时保持可靠。实证表明，在线迭代对齐方法能将奖励过优化程度降低5至15个百分点\cite{tajwar2024preference,online_rlhf_2025}。

图~\ref{fig:reward_overopt_loop}概括了固定奖励模型下分布偏移与奖励过优化的形成链条。奖励模型在偏好数据上训练完成后被冻结，而策略仍在PPO或GRPO中持续更新；当策略生成分布逐渐进入奖励模型训练分布之外的区域时，评分误差会被策略优化反复放大，最终表现为代理奖励升高但真实质量下降。

\begin{figure}[htbp]
    \centering
    \resizebox{0.94\textwidth}{!}{%
    \begin{tikzpicture}[
        font=\scriptsize,
        stage/.style={draw, rounded corners, thick, fill=gray!7, minimum width=10.8cm, minimum height=4.15cm},
        box/.style={draw, rounded corners, thick, align=center, minimum width=2.45cm, minimum height=0.82cm, text width=2.45cm},
        warn/.style={draw, rounded corners, thick, align=center, minimum width=5.9cm, minimum height=0.76cm, text width=5.9cm},
        arrow/.style={->, thick}
    ]
        \node[stage] at (4.9,-0.22) {};
        \node[box, fill=yellow!15] (pref) at (1.45,0.95) {偏好数据\\训练分布};
        \node[box, fill=yellow!15] (rmtrain) at (4.9,0.95) {训练奖励模型\\$r_\psi$};
        \node[box, fill=yellow!15] (freeze) at (8.35,0.95) {奖励模型\\参数冻结};
        \node[box, fill=blue!8] (policy) at (8.35,-0.35) {策略持续优化\\$\pi_\phi$更新};
        \node[box, fill=orange!10] (shift) at (4.9,-0.35) {生成分布偏移\\分布外样本};
        \node[box, fill=red!8] (error) at (1.45,-0.35) {评分误差放大\\奖励过优化};
        \draw[arrow] (pref) -- (rmtrain);
        \draw[arrow] (rmtrain) -- (freeze);
        \draw[arrow] (freeze) -- (policy);
        \draw[arrow] (policy) -- (shift);
        \draw[arrow] (shift) -- (error);

        \node[warn, fill=red!8] (result) at (4.9,-1.65) {结果：代理奖励继续升高，真实回答质量可能下降};
        \draw[arrow] (error.south) |- (result.west);
    \end{tikzpicture}}
    \caption{固定奖励模型下的分布偏移与奖励过优化机制}
    \label{fig:reward_overopt_loop}
\end{figure}

\section{直接偏好优化}
\label{sec:dpo}

\subsection{从RLHF目标到DPO损失的推导}

直接偏好优化（Direct Preference Optimization, DPO）\cite{rafailov2023direct}从式~\eqref{eq:kl_constrained}的KL约束强化学习目标出发，通过代数变换把RLHF的两阶段流程（奖励建模与在线策略优化）压缩为对偏好数据的单阶段监督学习，省去了显式奖励模型和在线采样。

图~\ref{fig:dpo_pipeline}给出了DPO相对于RLHF的简化路径。DPO不再显式训练奖励模型，也不需要在训练过程中执行PPO/GRPO式在线策略优化，而是用参考模型与当前策略在偏好响应和非偏好响应上的对数概率比，直接构造偏好分类损失并更新策略参数。

\begin{figure}[htbp]
    \centering
    \resizebox{0.98\textwidth}{!}{%
    \begin{tikzpicture}[
        font=\normalsize,
        box/.style={draw, rounded corners, thick, align=center, minimum width=2.6cm, minimum height=0.95cm, text width=2.6cm},
        wide/.style={draw, rounded corners, thick, align=center, minimum width=3.2cm, minimum height=0.95cm, text width=3.2cm},
        arrow/.style={->, thick}
    ]
        \node[wide, fill=orange!10] (data) at (0,0) {偏好数据\\$(x,y_w,y_l)$};
        \node[wide, fill=blue!8] (prob) at (4.0,0) {当前/参考策略\\计算响应概率};
        \node[wide, fill=green!10] (logratio) at (8.0,0) {计算两类响应\\对数概率比};
        \node[wide, fill=yellow!15] (loss) at (12.0,0) {DPO偏好损失\\提升$y_w$、压低$y_l$};
        \node[wide, fill=purple!8] (update) at (16.0,0) {直接更新\\策略模型};
        \draw[arrow] (data) -- (prob);
        \draw[arrow] (prob) -- (logratio);
        \draw[arrow] (logratio) -- (loss);
        \draw[arrow] (loss) -- (update);
        \node[draw, dashed, rounded corners, align=center, text width=5.8cm, minimum height=0.9cm] at (8.0,-1.65)
            {被省略的RLHF组件：\\显式奖励模型训练、在线采样、价值/优势估计};
    \end{tikzpicture}}
    \caption{直接偏好优化的单阶段训练路径}
    \label{fig:dpo_pipeline}
\end{figure}

DPO推导的关键步骤如下。由式~\eqref{eq:boltzmann}，KL约束RL目标的最优策略满足：
\begin{equation}
    \pi^*(y|x) = \frac{\pi_{\text{ref}}(y|x) \exp(r^*(x,y)/\beta)}{Z(x)}
\end{equation}
对两边取对数并整理，可得奖励函数的策略参数化表示（即式~\eqref{eq:reward_reparam}）：
\begin{equation}
    r^*(x, y) = \beta \log \frac{\pi^*(y|x)}{\pi_{\text{ref}}(y|x)} + \beta \log Z(x)
    \label{eq:dpo_reparam}
\end{equation}
该式说明，给定最优策略 $\pi^*$，奖励函数 $r^*$ 可被唯一确定（至多相差一个仅依赖 $x$ 的加性常数 $\beta \log Z(x)$）。

将式~\eqref{eq:dpo_reparam}代入Bradley-Terry偏好概率式~\eqref{eq:bt_model}：由于 $\beta \log Z(x)$ 在两条响应的奖励差中相互抵消，可得：
\begin{align}
    \mathbb{P}^*(y_w \succ y_l | x) &= \sigma\left(r^*(x,y_w) - r^*(x,y_l)\right) \notag \\
    &= \sigma\left(\beta \log \frac{\pi^*(y_w|x)}{\pi_{\text{ref}}(y_w|x)} - \beta \log \frac{\pi^*(y_l|x)}{\pi_{\text{ref}}(y_l|x)}\right)
    \label{eq:dpo_preference}
\end{align}
式~\eqref{eq:dpo_preference}直接建立了偏好概率与策略参数之间的联系，无需显式奖励模型作为中介。以参数化策略 $\pi_\phi$ 替换最优策略 $\pi^*$，对偏好数据做最大似然估计，即得DPO损失：
\begin{equation}
    \mathcal{L}_{\text{DPO}}(\phi) = -\mathbb{E}_{(x, y_w, y_l) \sim \mathcal{D}}\left[\log \sigma\left(\beta \log \frac{\pi_\phi(y_w|x)}{\pi_{\text{ref}}(y_w|x)} - \beta \log \frac{\pi_\phi(y_l|x)}{\pi_{\text{ref}}(y_l|x)}\right)\right]
    \label{eq:dpo_loss}
\end{equation}
DPO把奖励模型隐式参数化于语言模型自身：给定策略 $\pi_\phi$，其隐式奖励为 $\hat{r}_\phi(x,y) = \beta \log \pi_\phi(y|x)/\pi_{\text{ref}}(y|x)$。最小化 $\mathcal{L}_{\text{DPO}}$ 等价于对该隐式奖励做偏好对比训练，使偏好响应的隐式奖励高于非偏好响应。

从信息论角度看，DPO损失可以解释为最小化参数化偏好分布 $\mathbb{P}_\phi(y_w \succ y_l | x)$ 与经验偏好分布之间的KL散度，等价于在偏好数据上做最大似然估计。DPO的训练过程可以看作在偏好数据对上做对比学习（contrastive learning），不需要任何额外的强化学习基础设施。其实现复杂度大幅低于PPO-RLHF，所需计算资源和工程投入也明显减少，这是DPO快速获得学术界和工业界关注的重要原因\cite{rafailov2023direct}。

\subsection{DPO的隐式奖励与梯度分析}

DPO虽不显式训练奖励模型，但损失函数中仍隐含了一个"奖励"量。由式~\eqref{eq:dpo_loss}可知，DPO实际比较的是当前策略与参考策略在同一响应上的生成概率比。下式可看作DPO对应的隐式奖励：
\begin{equation}
    \hat{r}_\phi(x, y) = \beta \log \frac{\pi_\phi(y|x)}{\pi_{\text{ref}}(y|x)} = \beta \sum_{t=1}^{|y|} \log \frac{\pi_\phi(y_t | x, y_{<t})}{\pi_{\text{ref}}(y_t | x, y_{<t})}
    \label{eq:dpo_implicit_reward}
\end{equation}
这里的"奖励"不是单独训练得到的奖励模型分数，而是当前策略相对参考策略生成某个响应的倾向。该值为正，说明当前策略比参考策略更倾向于生成该响应；为负则说明当前策略相对更不倾向于生成该响应。DPO的关键不是直接判断回答的绝对质量，而是调整当前策略相对参考策略的偏好方向。

对于一个偏好数据对 $(x,y_w,y_l)$，DPO希望偏好响应 $y_w$ 的隐式奖励高于非偏好响应 $y_l$。记
\begin{equation}
    h_w = \beta \log \frac{\pi_\phi(y_w|x)}{\pi_{\text{ref}}(y_w|x)}, \qquad
    h_l = \beta \log \frac{\pi_\phi(y_l|x)}{\pi_{\text{ref}}(y_l|x)} ,
\end{equation}
则DPO损失可以理解为推动 $h_w-h_l$ 变大。对策略参数 $\phi$ 求梯度可得：
\begin{align}
    \nabla_\phi \mathcal{L}_{\text{DPO}} = -\mathbb{E}_{(x,y_w,y_l) \sim \mathcal{D}} \Big[ &\sigma(h_l - h_w) \cdot \beta \Big(\nabla_\phi \log \pi_\phi(y_w|x) \notag \\
    &- \nabla_\phi \log \pi_\phi(y_l|x)\Big) \Big]
    \label{eq:dpo_gradient}
\end{align}
式~\eqref{eq:dpo_gradient}说明DPO的更新方向由两部分组成：提高偏好响应 $y_w$ 的生成概率，降低非偏好响应 $y_l$ 的生成概率。前面的权重 $\sigma(h_l-h_w)$ 反映当前模型仍难区分该偏好对的程度：模型已经明显更偏好 $y_w$ 时，$h_w \gg h_l$，权重接近0，样本对梯度贡献很小；模型仍难以区分二者时权重较大，该样本会产生更强的更新。DPO会自动把更多训练强度分配给当前模型尚未学会区分的偏好对。

需要注意的是，DPO优化的是偏好响应相对非偏好响应、当前策略相对参考策略的概率比，而非偏好响应的绝对生成概率。换言之，只要 $y_w$ 相对 $y_l$ 的优势扩大，损失就可能下降；这并不必然意味着模型一定提高了 $y_w$ 的绝对概率。在离线偏好数据与当前策略分布差异较大时，这种相对概率优化可能导致似然位移（likelihood displacement）现象——模型满足了偏好对上的相对排序，却未必提升了真实生成质量\cite{azar2024general}。这也是DPO相对在线RLHF方法的重要局限之一。

\subsection{DPO的变体方法}

DPO发表后，学术界从多个角度提出了若干改进变体。\textbf{身份偏好优化（IPO）}\cite{azar2024general}把DPO的对数似然目标替换为对奖励差值的二次惩罚，要求 $h_w - h_l = 1/(2\beta)$ 维持在固定值而非趋向无穷大，缓解了数据量有限时的过拟合倾向。\textbf{卡尼曼-特沃斯基优化（KTO）}\cite{ethayarajh2024kto}受前景理论启发，放弃成对偏好假设，改为对单条响应的绝对质量直接建模，仅需单点标注就能完成对齐，降低了数据收集成本。\textbf{简单偏好优化（SimPO）}\cite{meng2024simpo}以序列平均对数概率代替对数概率比作为奖励代理，去除了DPO中对响应长度的系统性偏好，并引入目标奖励差 $\gamma$ 以避免优化结果过于接近决策边界；该方法不需要参考模型前向推断，显存效率更高。\textbf{信赖域DPO（TR-DPO）}\cite{gorbatovski2025trdpo}在训练中动态刷新参考策略，减少对分布外数据的依赖，在对话与摘要基准上优于标准DPO。\textbf{词元级DPO（TDPO）}\cite{zeng2024token}把偏好信号从序列级扩展到词元级，通过细粒度信用分配提供更精确的训练梯度。

\subsection{DPO的优势与局限}

DPO的主要优势体现在几个方面。\textbf{训练简洁性}：DPO把三阶段RLHF流程压缩为单阶段对比学习，不必维护价值函数网络、设计奖励塑形（reward shaping）策略或调节PPO的众多超参数，实现成本明显降低。\textbf{理论等价性}：在Bradley-Terry假设和无限数据极限下，DPO与RLHF具有相同的最优策略，理论上可以恢复式~\eqref{eq:kl_constrained}的全局最优解。\textbf{训练稳定性}：DPO损失函数连续可微，梯度方差远小于基于采样的策略梯度方法，训练过程更平稳，不会出现PPO中常见的奖励崩溃（reward collapse）问题。

DPO存在若干固有局限，主要源于其离线（offline）学习的本质。\textbf{分布偏移问题}：DPO在固定偏好数据集 $\mathcal{D}$ 上训练，数据集中的响应来自数据收集时的策略（通常是SFT模型或其他固定模型），而不是当前训练中的策略 $\pi_\phi$。$\pi_\phi$ 参数不断更新，其生成分布与 $\mathcal{D}$ 中的数据分布会出现明显偏差，训练信号也就逐渐脱离当前策略的实际分布，影响最终对齐效果\cite{tajwar2024preference}。研究表明，使用在线采样的偏好数据（即从当前策略采样响应后收集偏好）明显优于离线固定数据集上的DPO训练，印证了在线数据的重要性。

\textbf{奖励信号静态性}：DPO把奖励隐式嵌入策略参数，缺乏对奖励信号的显式控制能力。RLHF中奖励模型可作为独立模块被重新评估、调整或替换；DPO中"奖励"只体现在策略对数概率比中，训练过程中难以单独调节其尺度或方向，也难以解释模型究竟学到了怎样的奖励偏好。\textbf{偏好标注噪声敏感性}：BT假设默认偏好比较能形成相对一致的排序；如果不同标注者的选择互相矛盾，或偏好强弱差异很大，DPO得到的训练信号就会变得不稳定。实验表明，偏好数据中噪声标注较多时DPO性能明显下降，而带鲁棒性约束的方法（如IPO）表现更稳健\cite{azar2024general}。

\textbf{长度偏差}：DPO以总对数概率（而非平均对数概率）计算隐式奖励，模型由此倾向于生成更长的响应——更长的序列在绝对概率比上更容易取得大值，即便生成质量并无实质提升\cite{meng2024simpo}。这一长度偏差在实际部署中可能导致模型输出冗余，需要额外的后处理或长度惩罚机制加以修正。

DPO用更简单的离线训练替代了RLHF中的显式奖励模型和在线策略优化，工程复杂度明显下降；但它也把奖励信号隐藏进策略概率比中，奖励难以被单独检查、更新和约束。当训练目标重新回到“奖励信号如何获得、如何随策略变化而调整”这一问题时，就需要引入从行为中反推奖励函数的逆强化学习框架。


\section{逆强化学习与奖励函数学习}

RLHF保留了显式奖励模型，但奖励模型通常是静态的；DPO简化了训练流程，但奖励信号变为隐式形式。若希望奖励信号能随策略分布变化而更新，就需要回到奖励函数学习本身。逆强化学习就是从专家行为中反推奖励函数的代表性框架。

\subsection{问题形式化与研究动机}

逆强化学习（Inverse Reinforcement Learning, IRL）\cite{ng2000algorithms}是强化学习的逆向问题。正向RL在已知奖励函数的MDP中寻找最优策略；IRL从专家（expert）在未知奖励函数下表现出的最优或近优行为轨迹中，反向推断潜在奖励函数 $r^*$，再以该奖励函数驱动正向RL进行策略学习。IRL的基本假设是：专家行为在某一潜在奖励函数下（近似）最优，因此观测到的专家轨迹中蕴含了关于任务目标的信息。

图~\ref{fig:rl_irl_relation}从信息流向角度概括了正向强化学习与逆强化学习的区别：正向RL以奖励函数为起点，通过策略优化得到行为；IRL从专家行为出发，先反推能解释该行为的奖励函数，再用该奖励函数训练策略。在大语言模型偏好对齐中，专家响应可看作偏好目标的行为体现，IRL则提供了一种从高质量示例中恢复偏好奖励的学习框架。这一视角使“偏好对齐”不再只依赖固定偏好分类器或静态奖励模型，而可以转化为随策略分布共同演化的奖励函数学习问题。

\begin{figure}[htbp]
    \centering
    \resizebox{0.98\textwidth}{!}{%
    \begin{tikzpicture}[
        font=\normalsize,
        box/.style={draw, rounded corners, thick, align=center, minimum width=2.65cm, minimum height=0.95cm, text width=2.65cm},
        wide/.style={draw, rounded corners, thick, align=center, minimum width=3.05cm, minimum height=0.95cm, text width=3.05cm},
        arrow/.style={->, thick}
    ]
        \node[anchor=east, font=\normalsize\bfseries] at (0,1.25) {正向RL};
        \node[wide, fill=blue!8] (rl_reward) at (2.2,1.25) {已知奖励函数\\$r(s,a)$};
        \node[wide, fill=blue!8] (rl_opt) at (6.0,1.25) {正向RL\\策略优化};
        \node[wide, fill=blue!8] (rl_policy) at (9.8,1.25) {最优/近优策略\\$\pi^*$};
        \node[wide, fill=blue!8] (rl_behavior) at (13.6,1.25) {策略行为轨迹\\$\tau$};
        \draw[arrow] (rl_reward) -- (rl_opt);
        \draw[arrow] (rl_opt) -- (rl_policy);
        \draw[arrow] (rl_policy) -- (rl_behavior);

        \node[anchor=east, font=\normalsize\bfseries] at (0,-1.25) {逆强化学习};
        \node[wide, fill=orange!10] (irl_demo) at (2.2,-1.25) {专家示例轨迹\\$\mathcal{D}_E$};
        \node[wide, fill=orange!10] (irl_infer) at (6.0,-1.25) {逆向推断\\奖励函数};
        \node[wide, fill=orange!10] (irl_reward) at (9.8,-1.25) {学习奖励\\$r_\theta$};
        \node[wide, fill=orange!10] (irl_policy) at (13.6,-1.25) {策略优化并\\逼近专家行为};
        \draw[arrow] (irl_demo) -- (irl_infer);
        \draw[arrow] (irl_infer) -- (irl_reward);
        \draw[arrow] (irl_reward) -- (irl_policy);
    \end{tikzpicture}}
    \caption{正向强化学习与逆强化学习的信息流对比}
    \label{fig:rl_irl_relation}
\end{figure}

Ng和Russell的奠基性工作\cite{ng2000algorithms}给出了IRL问题的形式化与可行性分析。给定MDP $\mathcal{M} \setminus \mathcal{R}$（即除奖励函数外的所有要素）和专家策略 $\pi_E$（或专家轨迹集合 $\mathcal{D}_E = \{\tau_E^{(i)}\}$），IRL的目标是恢复奖励函数 $r^*$，使专家行为在该奖励下具有较高回报并接近最优。但IRL存在固有的不适定性（ill-posedness）：能解释同一组专家行为的奖励函数通常不止一个。例如把所有状态和动作的奖励都设为零时，任意策略都同样"最优"，但这样的奖励函数显然无法揭示专家行为背后的偏好。IRL方法需要额外的准则，从多个候选奖励函数中选出更合理的解。

IRL与模仿学习（Imitation Learning, IL）密切相关，但二者有本质区别\cite{li2024imitation}。行为克隆（Behavior Cloning, BC）直接从专家示例中学习状态-动作映射 $\pi_E: s \mapsto a$，归结为监督学习问题，但面临复合误差（compounding error）困境：在训练分布之外，微小的预测误差可能不断累积并引发分布偏移。IRL通过显式建模奖励函数，所学策略具备更强的泛化能力，也能处理状态分布偏移的场景。在大语言模型对齐中，IRL的优势在于：所学奖励模型 $r_\theta$ 可以随当前策略分布变化而更新，缓解固定奖励模型面对的分布偏移问题。这也带来新的要求：奖励函数不能在多轮更新中任意漂移，否则策略优化会受不稳定奖励信号的影响。

IRL也可以从最优控制与概率推断的视角加以理解。在控制即推断（Control as Inference）框架下\cite{ziebart2008maximum}，轨迹的最优性被视为潜变量，IRL等价于在给定专家轨迹观测时对奖励函数做后验推断。这一概率视角为最大熵IRL提供了理论基础，也展现了IRL与变分推断、能量模型等概率图模型领域之间的联系。

\subsection{最大熵逆强化学习}

最大熵逆强化学习（Maximum Entropy IRL, MaxEnt IRL）由Ziebart等人\cite{ziebart2008maximum}提出，为上述不适定性问题给出了一个概率化解决思路。其基本假设是：专家轨迹之所以被观察到，是因为它们在某个潜在奖励函数下具有较高概率；对数据没有明确区分的轨迹，则不应人为赋予过强偏好。最大熵原理用来表达这一点：在解释专家示例的同时，尽量保留轨迹选择的不确定性。MaxEnt IRL不仅寻找一个使专家策略最优的奖励函数，而且同时学习一个由奖励函数诱导的轨迹概率分布。

在上述原则下，MaxEnt IRL要为每一条可能轨迹分配一个概率。记 $p_\theta(\tau)$ 为奖励参数 $\theta$ 下轨迹 $\tau$ 被生成的概率。一条轨迹获得的累积奖励越高，它在模型分布下应越可能出现；最大熵推导给出的具体形式为：
\begin{equation}
    p_\theta(\tau) = \frac{\exp\left(r_\theta(\tau)\right)}{Z(\theta)} = \frac{\exp\left(\sum_t r_\theta(s_t, a_t)\right)}{Z(\theta)}
    \label{eq:maxent_distribution}
\end{equation}
其中，$r_\theta(\tau)=\sum_t r_\theta(s_t,a_t)$ 表示整条轨迹的累积奖励；$Z(\theta)=\sum_\tau \exp(r_\theta(\tau))$ 为归一化配分函数，保证所有轨迹的概率之和为1。式~\eqref{eq:maxent_distribution}说明，MaxEnt IRL把奖励函数转化为轨迹上的概率模型：高奖励轨迹具有更高概率，但低奖励轨迹也不会被完全排除，这与最大熵原理保留不确定性的思想一致。

基于该轨迹分布，MaxEnt IRL通过提高专家轨迹在模型分布下的概率来学习奖励函数。给定专家轨迹数据集 $\mathcal{D}_E = \{\tau_E^{(i)}\}_{i=1}^N$，其优化目标为最大化专家轨迹的平均对数似然：
\begin{equation}
    \max_\theta\, \mathcal{L}_{\text{MaxEnt}}(\theta) = \frac{1}{N}\sum_{i=1}^N \log p_\theta(\tau_E^{(i)}) = \frac{1}{N}\sum_{i=1}^N r_\theta(\tau_E^{(i)}) - \log Z(\theta)
    \label{eq:maxent_objective}
\end{equation}

对式~\eqref{eq:maxent_objective}的直观理解是：第一项提高专家轨迹的奖励，使专家示例在模型分布下更可能出现；第二项 $\log Z(\theta)$ 由所有可能轨迹共同决定，避免模型仅靠无界增大奖励值来提高似然。对该目标求梯度，可得到“提高专家轨迹奖励、降低模型自身高概率轨迹奖励”的更新方向。换言之，MaxEnt IRL通过对比专家轨迹与当前模型轨迹来学习奖励函数：专家轨迹获得更高奖励，模型当前容易生成但与专家行为不同的轨迹则被压低。

大规模状态动作空间中，配分函数 $Z(\theta)$ 需要对所有可能轨迹求和，通常无法直接计算。经典MaxEnt IRL可借助动态规划估计轨迹分布；复杂模型和连续空间中则常通过采样或对比学习近似。

MaxEnt IRL与能量模型（Energy-Based Model, EBM）之间存在紧密联系\cite{wulfmeier2024irl}。把 $-r_\theta(\tau)$ 看作轨迹的能量，式~\eqref{eq:maxent_distribution}就是以轨迹为随机变量的能量模型。训练能量模型的标准方法对比散度（Contrastive Divergence）与上述最大似然目标的梯度结构一致：增大数据样本（专家轨迹）的概率，同时减小模型生成样本（当前策略轨迹）的概率。配分函数 $Z(\theta)$ 在连续空间中通常无法解析计算，这是MaxEnt IRL的主要计算难点，也推动了后续基于MCMC的采样方法和对比方法（如GAIL）的发展。

在语言模型对齐的序列级奖励设定下，MaxEnt IRL目标简化为如下形式：
\begin{align}
    \mathcal{L}_{\text{IRL}}(\theta) &= \mathbb{E}_{(x,y) \sim \mathcal{D}_E}\left[r_\theta(x, y)\right] - \log Z(\theta) \notag \\
    \nabla_\theta \mathcal{L}_{\text{IRL}} &= \mathbb{E}_{(x,y_E) \sim \mathcal{D}_E}\left[\nabla_\theta r_\theta(x, y_E)\right] - \mathbb{E}_{x \sim \mathcal{D}, y \sim \pi^*_\theta(\cdot|x)}\left[\nabla_\theta r_\theta(x, y)\right]
    \label{eq:maxent_seq}
\end{align}
换言之，奖励学习通过专家样本与模型当前样本之间的对比来更新：专家响应应获得更高解释概率，模型当前容易生成但不符合专家分布的响应则相对被压低。这一对比形式也是本文第三章方法设计的理论基础。

\subsection{生成对抗模仿学习}

生成对抗模仿学习（Generative Adversarial Imitation Learning, GAIL）把IRL与生成对抗网络（GAN）结合：奖励函数类比为判别器 $D_\theta$，策略类比为生成器，通过 $\min_{\pi_\phi} \max_{D_\theta}$ 对抗训练交替更新，绕开了MaxEnt IRL中配分函数的显式计算。从信息论角度，GAIL最小化当前策略与专家占用度量之间的Jensen-Shannon散度；判别器无需手工特征设计，由深度网络端到端学习。


\subsection{奖励—策略双层优化框架}

把IRL框架应用于语言模型对齐时，本文采用序列级单步决策表述：给定提示 $x$ 后，完整响应 $y$ 看作一次动作，动作价值函数等价写为序列级奖励 $Q(x,y)=r_\theta(x,y)$。在该设定下，奖励学习不是先固定一个奖励模型再单独训练策略，而是通过奖励诱导出的最优策略来解释专家响应。原始问题可写为：
\begin{equation}
\begin{aligned}
    \max_{\theta}\quad \ell(\theta)
    &:= \mathbb{E}_{x \sim \rho,\, y \sim \pi_{\mathrm{E}}(\cdot|x)}
    \left[\log \pi_\theta(y|x)\right], \\
    \text{s.t.}\quad
    \pi_\theta
    &:= \arg\max_{\pi}\ \mathbb{E}_{x \sim \rho}\left[
    \mathbb{E}_{y \sim \pi(\cdot|x)}r_\theta(x,y)
    -\beta D_{\mathrm{KL}}\!\left(\pi(\cdot|x)\|\pi_{\mathrm{ref}}(\cdot|x)\right)\right].
    \label{eq:bilevel_original_ch2}
\end{aligned}
\end{equation}
其中，$\pi_\theta$ 表示由奖励函数 $r_\theta$ 通过内层KL正则策略优化诱导出的最优策略，$\pi_{\mathrm{E}}$ 表示专家响应分布，$\pi_{\mathrm{ref}}$ 为参考策略，$\beta$ 控制策略偏离参考模型的程度。直观来看，外层目标负责学习奖励函数：调整奖励参数 $\theta$，使该奖励函数诱导出的策略更倾向于生成专家响应；内层目标在给定奖励函数下优化生成策略，由KL正则约束新策略相对参考策略的偏离。

内层问题具有闭式解。对固定提示 $x$，由KKT条件可得：
\begin{equation}
    \pi_\theta(y|x)=
    \frac{\pi_{\mathrm{ref}}(y|x)\exp\left(r_\theta(x,y)/\beta\right)}
    {\sum_{\tilde{y}\in\mathcal{Y}}\pi_{\mathrm{ref}}(\tilde{y}|x)\exp\left(r_\theta(x,\tilde{y})/\beta\right)}.
    \label{eq:bilevel_policy_ch2}
\end{equation}
该式说明，奖励函数并不直接等同于最终策略，而是通过对参考策略做指数加权来改变响应分布；高奖励响应的概率被提高，KL项则约束这种改变的幅度。

把式~\eqref{eq:bilevel_policy_ch2}代回外层目标，并令
\begin{equation}
    J(\pi,r_\theta)
    :=\mathbb{E}_{x\sim\rho}\left[
    \mathbb{E}_{y\sim\pi(\cdot|x)}r_\theta(x,y)
    -\beta D_{\mathrm{KL}}\!\left(\pi(\cdot|x)\|\pi_{\mathrm{ref}}(\cdot|x)\right)
    \right],
    \label{eq:regularized_j_ch2}
\end{equation}
则在忽略与 $\theta$ 无关的专家策略KL常数 $-\beta D_{\mathrm{KL}}(\pi_{\mathrm{E}}\|\pi_{\mathrm{ref}})$ 后，可得到等价的奖励学习目标：
\begin{equation}
    \max_\theta\quad L(\theta)=J(\pi_{\mathrm{E}},r_\theta)-J(\pi_\theta,r_\theta).
    \label{eq:bilevel_outer_ch2}
\end{equation}
需要强调的是，式~\eqref{eq:bilevel_outer_ch2}是最大化问题，不是最小化"专家奖励与当前策略奖励的差距"。$\pi_\theta$ 是当前奖励下的KL正则最优策略，$J(\pi_\theta,r_\theta)$ 对应归一化项或软最优值；最大化 $J(\pi_{\mathrm{E}},r_\theta)-J(\pi_\theta,r_\theta)$ 的含义是提高专家响应在由奖励诱导的策略分布中的概率，使专家策略在当前奖励下尽量接近软最优，而不是简单地把专家样本和模型样本的裸奖励差写成一个任意方向的分类目标。

实际训练中，$\pi_\theta$ 随 $\theta$ 改变而改变，精确求解内层最优策略代价较高。可在第 $k$ 轮使用旧策略 $\pi_{\mathrm{old}}$ 构造局部替代目标：
\begin{equation}
    L_{\mathrm{old}}(\theta)=J(\pi_{\mathrm{E}},r_\theta)-J(\pi_{\mathrm{old}},r_\theta).
    \label{eq:bilevel_surrogate_ch2}
\end{equation}
当 $\pi_{\mathrm{old}}=\pi_{\theta_{\mathrm{old}}}$ 时，该替代目标与原始目标在 $\theta_{\mathrm{old}}$ 处具有相同的函数值和一阶导数，可作为原始双层目标的局部近似，用于交替更新奖励模型与策略模型。这一双层目标说明，IRL式对齐不只是训练一个奖励模型，而是在奖励学习和策略优化之间反复交替。由此带来的关键问题是：奖励函数需要更新，但更新幅度又要受到控制。第三章会在该目标基础上分析奖励更新幅度对单调改进下界的影响，并据此设计约束调节机制。


\section{各对齐方法综合对比}

\subsection{方法体系梳理}

前文按照策略优化、奖励建模、偏好优化和奖励函数学习的顺序介绍了相关方法。PPO和GRPO主要解决“给定奖励后如何更新策略”；RLHF进一步回答“奖励如何由人类偏好数据训练得到”；DPO则尝试绕开奖励模型和在线强化学习，直接用偏好数据更新策略；IRL重新回到奖励函数学习本身，关注如何从专家示例中恢复可用于策略优化的奖励信号。

从奖励信号来看，RLHF采用显式奖励模型，奖励分数可以直接观察和分析，但奖励模型通常在策略优化前固定，后续策略分布变化会削弱评分可靠性。DPO及其变体不再单独训练奖励模型，而是把奖励信号隐含在策略与参考策略的概率比中，因此训练流程更短，但也难以单独检查、调整和约束奖励。IRL和GAIL从专家示例中学习奖励或判别信号，并在训练过程中让奖励与策略相互作用，更适合处理策略分布持续变化的情形，但训练过程更复杂。

从优化方式来看，PPO依赖价值函数估计优势，训练链路较重；GRPO用组内相对奖励替代价值函数估计，降低了工程成本，但两类方法仍依赖外部奖励信号。RLHF把奖励模型与PPO/GRPO结合，形成显式奖励驱动的在线对齐流程；DPO、IPO、KTO、SimPO等方法主要依赖固定离线偏好数据，以监督学习形式完成训练；IRL和GAIL则需要在线采样或交替优化，使奖励学习和策略更新形成闭环。这些方法的差异主要体现在训练简洁性、奖励可控性和分布适应性之间的取舍。

\subsection{综合对比分析}

表~\ref{tab:method_comparison}对现有主要对齐方法做对比，重点关注优化方式、数据类型、计算开销、分布适应性和工程复杂度。其中分布适应性指当前策略生成的内容逐渐偏离原始训练数据时，方法是否具有调整奖励信号或训练信号的机制。

\begin{table}[htbp]
    \centering
    \caption{大语言模型对齐方法综合对比}
    \label{tab:method_comparison}
    \small
    \setlength{\tabcolsep}{10pt}
    \renewcommand{\arraystretch}{1.3}
    \begin{tabular}{lcccc c}
        \toprule
        \textbf{方法} & \textbf{优化方式} & \textbf{数据类型} & \textbf{计算开销} & \textbf{分布适应性} & \textbf{工程复杂度} \\
        \midrule
        RLHF/PPO   & 在线RL   & 偏好对   & 高 & 低 & 高 \\
        RLHF/GRPO  & 在线RL   & 偏好对   & 中 & 低 & 中 \\
        DPO        & 离线对比 & 偏好对   & 低 & 低 & 低 \\
        IPO        & 离线对比 & 偏好对   & 低 & 中 & 低 \\
        KTO        & 离线对比 & 单点标注 & 低 & 中 & 低 \\
        SimPO      & 离线对比 & 偏好对   & 低 & 中 & 低 \\
        MaxEnt IRL & 在线RL   & 专家示例 & 高 & 高 & 高 \\
        GAIL       & 在线RL   & 专家示例 & 高 & 高 & 高 \\
        \bottomrule
    \end{tabular}
\end{table}

由表~\ref{tab:method_comparison}可以看出，现有方法大致形成三类路线。第一类是显式奖励驱动的RLHF。PPO版本保留价值函数和优势估计，训练链路完整但成本较高；GRPO版本省去独立价值函数，工程上更轻量，但两种方式都依赖固定奖励模型，在策略持续更新后容易遇到分布偏移和奖励模型过优化。第二类是DPO及其变体。这类方法把偏好学习转化为离线对比训练，不必单独维护奖励模型和在线采样流程，更容易实现和调参；但它们主要依赖固定数据集，缺少根据当前策略输出主动更新奖励信号的机制。第三类是IRL和GAIL。这类方法通过专家示例学习奖励或判别信号，允许奖励信号与策略更新相互影响，更适合处理策略分布持续变化时的对齐问题；这种优势也伴随着更高的采样成本和稳定性要求。

现有对齐方法没有单一意义上的最优选择。PPO和GRPO改善的是策略优化环节，并不能保证奖励信号本身可靠；DPO简化了训练流程，但牺牲了显式奖励的可检查性和可更新性；IRL能够动态学习奖励，但要应对奖励—策略交替更新带来的稳定性问题。由此可见，本文所称“基于逆强化学习”并不是简单地用IRL替代RLHF或DPO，而是把偏好对齐的核心问题重新表述为动态奖励函数学习：奖励需要随策略变化不断更新，同时又必须避免更新幅度过大、训练过程振荡或奖励尺度漂移。这正是后续章节讨论动态奖励更新约束的直接背景。

\section{本章小结}

本章围绕大语言模型对齐中的策略优化与奖励信号构造展开，按PPO/GRPO、RLHF、DPO和IRL的顺序梳理了相关技术基础，为第三章的受约束逆强化学习方法做准备。

强化学习部分给出了基本框架以及PPO、GRPO两类策略优化算法。马尔可夫决策过程把语言模型生成看作逐步选择词元的决策过程；策略梯度方法说明了模型如何根据奖励信号调整生成概率。PPO通过截断概率比限制单次策略更新幅度，训练较稳定但需要价值函数；GRPO通过同一提示下多条响应的组内相对奖励估计优势，降低了价值函数带来的工程成本。这两类算法主要回答“如何根据奖励更新模型”。

在此基础上，本章讨论了RLHF的奖励建模与策略优化流程。RLHF用人类偏好数据训练奖励模型，再结合PPO或GRPO优化策略，是典型的显式奖励建模方法。其特点是奖励分数清晰、可直接用于策略更新；但奖励模型通常在策略优化前固定，当策略生成分布逐渐偏离奖励模型训练分布时，评分可靠性会下降，并可能进一步引发奖励模型过优化。

之后，本章梳理了DPO及其变体。DPO利用KL约束策略优化中的奖励—策略等价关系，把偏好学习转化为离线对比训练，省去了显式奖励模型和在线强化学习步骤，训练更简单稳定。但DPO的奖励信号隐含在策略与参考策略的概率比中，难以单独检查、调整和约束；固定离线数据也限制了它对当前策略分布变化的适应能力。

本章最后介绍了IRL、MaxEnt IRL和GAIL等奖励函数学习方法。IRL从专家示例中反推奖励函数，为动态学习奖励信号提供了理论框架；MaxEnt IRL用概率形式表达高奖励行为更可能出现的思想；GAIL从对抗学习角度实现专家示例学习。相较固定奖励模型或隐式奖励方法，IRL类方法能让奖励信号随训练过程更新，但也带来多轮奖励—策略交替更新的稳定性问题。如何在保持奖励函数可更新的同时控制其变化幅度，是下一章需要重点解决的问题。由此，相关技术基础最终汇聚到本文的核心问题：在大模型偏好对齐中，如何把IRL的奖励反演能力转化为稳定可训练的动态奖励优化方法。
